\clearpage
\subsection{Exercises}
\label{sec:parameter-estimation-exercises}

\subsubsection*{(A) James--Stein versus Least Squares}
\phantomsection
\addcontentsline{toc}{subsubsection}{(A) James--Stein versus Least Squares}
\label{ex:james-stein}

Least squares estimates each normal mean by its observed value. Can we improve on this natural rule
by shrinking the coordinates together, even when the true means are unrelated? The challenge is to
compare risks without knowing the means. We first construct an unbiased estimate of risk, then use it to
show that James--Stein estimation improves on least squares at every parameter value.

Let $X\sim\mathcal N_n(\theta,I_n)$, where $\theta\in\mathbb R^n$ is unknown. Measure the risk of an estimator by its total mean squared error,
\[
 R(\theta,\hat\theta)=\mathbb E_\theta\bigl[\|\hat\theta(X)-\theta\|_2^2\bigr].
\]
The least-squares estimator is $\hat\theta_{\mathrm{LS}}(X)=X$.

\begin{enumerate}[label=\textcolor{AccentInk}{\sffamily\bfseries A\arabic*.},ref=A\arabic*,leftmargin=2.65em,labelsep=0.65em,itemsep=1em]
\item \textbf{(Stein's unbiased risk estimate)} Let $\hat\theta:\mathbb R^n\to\mathbb R^n$ be weakly differentiable, with $\mathbb E_\theta\|\hat\theta(X)\|_2^2<\infty$ and $\mathbb E_\theta|\partial\hat\theta_i(X)/\partial x_i|<\infty$ for each $i$. Define
\[
 \operatorname{SURE}(X;\hat\theta)=\|\hat\theta(X)-X\|_2^2+2\sum_{i=1}^n\frac{\partial\hat\theta_i}{\partial x_i}(X)-n.
\]
Prove that
\[
 \mathbb E_\theta[\operatorname{SURE}(X;\hat\theta)]=\mathbb E_\theta\bigl[\|\hat\theta(X)-\theta\|_2^2\bigr].
\]
Why can SURE be computed from the data even though the loss cannot? \emph{Hint:} Expand the squared error and integrate by parts against the Gaussian density.

\item \textbf{(James--Stein dominates least squares)} For $n\ge3$, consider the James--Stein estimator
\[
 \hat\theta_{\mathrm{JS}}(X)=\left(1-\frac{n-2}{\|X\|_2^2}\right)X,\qquad X\ne0,
\]
with $\hat\theta_{\mathrm{JS}}(0)=0$. Calculate its SURE and use A1 to find the exact risk difference $R(\theta,\hat\theta_{\mathrm{JS}})-R(\theta,\hat\theta_{\mathrm{LS}})$. Deduce that, for every $\theta\in\mathbb R^n$,
\[
 R(\theta,\hat\theta_{\mathrm{JS}})<R(\theta,\hat\theta_{\mathrm{LS}})=n.
\]
In applying A1, check the regularity at the origin and show that $\mathbb E_\theta[\|X\|_2^{-2}]<\infty$. Where does the condition $n\ge3$ enter the argument?
\end{enumerate}

\subsubsection*{(B) Bayes Optimal Estimation for Gaussian Models}
\phantomsection
\addcontentsline{toc}{subsubsection}{(B) Bayes Optimal Estimation for Gaussian Models}
\label{ex:gaussian-bayes}

A Gaussian prior lets us combine observations with prior information about the unknown parameter.
The resulting estimator minimizes average risk under that prior. How does it compare with least squares,
and which estimator protects us against the worst possible parameter? This problem connects these two
notions of optimality by making the prior increasingly diffuse.

Under squared error loss, consider the Gaussian linear model and Gaussian prior
\[
 X\mid\theta\sim P_\theta=\mathcal N_n(A\theta,\sigma^2 I_n),\qquad
 \theta\sim Q=\mathcal N_p(0,\Sigma_0),
\]
where $A\in\mathbb R^{n\times p}$ is known and has full column rank, $\sigma^2>0$ is known, and $\Sigma_0\succ0$.

\begin{enumerate}[label=\textcolor{AccentInk}{\sffamily\bfseries B\arabic*.},ref=B\arabic*,leftmargin=2.65em,labelsep=0.65em,itemsep=1em]
\item \textbf{(Bayes optimal estimator)} Find the posterior distribution of $\theta$ given $X$. Use it to derive the Bayes estimator $\hat\theta_Q(X)$ and explain why it minimizes posterior expected squared error. Calculate its Bayes risk $\int_{\mathbb R^p}R(\theta,\hat\theta_Q)\,Q(\mathrm d\theta)$.

\item \textbf{(Least-squares estimator)} Find
\[
 \hat\theta_{\mathrm{LS}}(X)=\operatorname*{argmin}_{u\in\mathbb R^p}\|X-Au\|_2^2
\]
and calculate $R(\theta,\hat\theta_{\mathrm{LS}})$. Does this risk depend on $\theta$? Show that the estimator in B1 converges to $\hat\theta_{\mathrm{LS}}(X)$ when $\Sigma_0=\tau^2 I_p$ and $\tau\to\infty$.

\item \textbf{(Minimax optimality)} Determine the minimax risk
\[
 \inf_{\hat\theta}\sup_{\theta\in\mathbb R^p}R(\theta,\hat\theta)
\]
and identify an estimator that attains it. \emph{Hint:} The maximum risk is at least the average risk under any prior. Use B1 with $Q_\tau=\mathcal N_p(0,\tau^2 I_p)$ to obtain a lower bound, then let $\tau\to\infty$ and compare with B2.

Is the Bayes estimator for a fixed $\Sigma_0\succ0$ also minimax over $\mathbb R^p$? Finally, when $A=I_n$ and $\sigma=1$, explain why your minimax conclusion is compatible with the strict improvement in A2.
\end{enumerate}

\subsubsection*{(C) Tweedie's Theorem}
\phantomsection
\addcontentsline{toc}{subsubsection}{(C) Tweedie's Theorem}
\label{ex:tweedie}

Tweedie's formula recovers the posterior mean from the marginal density of Gaussian observations. A
second derivative reveals posterior uncertainty, even for a non-Gaussian prior~\cite{efron2011tweedie}.

Let $X\mid\theta\sim\mathcal N_p(\theta,\sigma^2 I_p)$, with known $\sigma^2>0$, and let $\theta\sim Q$ for any prior satisfying $\mathbb E_Q\|\theta\|_2^2<\infty$. Write the marginal density of $X$ as
\[
 m(x)=\int_{\mathbb R^p}\phi_\sigma(x-\theta)\,Q(\mathrm d\theta),\qquad
 \phi_\sigma(z)=\frac{\exp\{-\|z\|_2^2/(2\sigma^2)\}}{(2\pi\sigma^2)^{p/2}}.
\]
All gradients and Hessians below are with respect to $x$.

\begin{enumerate}[label=\textcolor{AccentInk}{\sffamily\bfseries C\arabic*.},ref=C\arabic*,leftmargin=2.65em,labelsep=0.65em,itemsep=1em]
\item \textbf{(Posterior mean from the marginal density)} Differentiate $m(x)$ under the integral sign and use Bayes' rule to prove
\[
 \hat\theta_Q(x)=\mathbb E[\theta\mid X=x]=x+\sigma^2\nabla\log m(x).
\]
This is Tweedie's formula. Why is it the Bayes optimal estimator under squared error loss? Interpret the direction of the correction $\sigma^2\nabla\log m(x)$.

\item \textbf{(Posterior covariance from curvature)} Differentiate once more and show that
\[
 \operatorname{Cov}(\theta\mid X=x)=\sigma^2 I_p+\sigma^4\nabla^2\log m(x).
\]
Deduce that $\nabla^2\log m(x)\succeq-\sigma^{-2}I_p$. If $m$ is log-concave, how does the posterior covariance compare with the noise covariance $\sigma^2 I_p$?

\item \textbf{(Returning to the Gaussian prior)} Take $Q=\mathcal N_p(0,\Sigma_0)$ with $\Sigma_0\succ0$. Find $m(x)$ and use C1--C2 to calculate the posterior mean and covariance. Verify that both agree with B1 when $A=I_p$.
\end{enumerate}

\subsubsection*{(D) Complete Statistics in the Exponential Model}
\phantomsection
\addcontentsline{toc}{subsubsection}{(D) Complete Statistics in the Exponential Model}
\label{ex:exponential-completeness}

Can the sum of observed lifetimes retain all the information about their mean? What does completeness
add to sufficiency? We will answer these questions and then use the sum to construct an optimal unbiased
estimator of the probability that a lifetime exceeds a fixed threshold.

Let $X_1,\ldots,X_n$ be independent with common density
\[
 f_\theta(x)=\frac1\theta e^{-x/\theta}\mathbf1_{\{x>0\}},\qquad\theta>0,
\]
where $n\ge2$ and $\theta$ is the unknown scale parameter. Write $T=\sum_{i=1}^nX_i$ and fix $t_0>0$.

\begin{enumerate}[label=\textcolor{AccentInk}{\sffamily\bfseries D\arabic*.},ref=D\arabic*,leftmargin=2.65em,labelsep=0.65em,itemsep=1em]
\item \textbf{(Sufficiency and the distribution of the sum)} Write the joint density and show that $T$ is sufficient for $\theta$. Use the likelihood-ratio characterization to prove that it is minimal sufficient. Find the density of $T$, including its normalizing constant.

\item \textbf{(Completeness)} Suppose a measurable function $g$ satisfies
\[
 \mathbb E_\theta|g(T)|<\infty,\qquad\mathbb E_\theta[g(T)]=0\quad\text{for every }\theta>0.
\]
Use the density from D1 to prove that $g(T)=0$ almost surely for every $\theta>0$. Conclude that $T$ is complete. \emph{Hint:} Set $s=1/\theta$ and use uniqueness of the Laplace transform.

\item \textbf{(Conditioning on the total)} Find the conditional density of $X_1$ given $T=t$. Hence calculate
\[
 \mathbb P_\theta(X_1>t_0\mid T=t)
\]
for every $t>0$, treating $t\le t_0$ and $t>t_0$ separately. Why does the conditional distribution not involve $\theta$?

\item \textbf{(From completeness to optimal unbiased estimation)} Show that $\mathbf1_{\{X_1>t_0\}}$ is unbiased for $e^{-t_0/\theta}$. Use D3 to write its Rao--Blackwell improvement explicitly as a function of $T$. Prove that the resulting estimator is the unique UMVUE of $e^{-t_0/\theta}$ and has strictly smaller variance than $\mathbf1_{\{X_1>t_0\}}$. Explain the different roles of sufficiency and completeness in this argument.
\end{enumerate}

\subsubsection*{(E) Pinsker's Linear Minimax Theorem}
\phantomsection
\addcontentsline{toc}{subsubsection}{(E) Pinsker's Linear Minimax Theorem}
\label{ex:pinsker-linear-minimax}

An ellipsoid gives different directions different signal budgets. How should a linear estimator shrink
each direction to minimize its worst-case risk? Following Lectures 16--17~\cite{wei2026stat9700}, we first construct an upper
bound, then turn a quadratic maximization into a Gaussian Bayes problem to obtain a matching lower
bound.

Observe $X=\theta+\sigma Z$, where $Z\sim\mathcal N_d(0,I_d)$ and $\sigma>0$ is known. For a known positive definite matrix $A$, define
\[
 \Theta(A)=\{\theta\in\mathbb R^d:\theta^\top A\theta\le1\},\qquad
 R_L(\Theta(A))=\inf_{M\in\mathbb R^{d\times d}}\sup_{\theta\in\Theta(A)}R(\theta,MX),
\]
where $R(\theta,\hat\theta)=\mathbb E_\theta\|\hat\theta-\theta\|_2^2$. The infimum is over all deterministic linear maps, including those that mix coordinates. Write $(u)_+=\max\{u,0\}$.

\begin{enumerate}[label=\textcolor{AccentInk}{\sffamily\bfseries E\arabic*.},ref=E\arabic*,leftmargin=2.65em,labelsep=0.65em,itemsep=1em]
\item \textbf{(Rotation and quadratic risk)} Show that an orthogonal change of coordinates preserves the model, the loss, and the class of linear estimators. Hence we may assume from now on that $A=\operatorname{diag}(a_1^2,\ldots,a_d^2)$ with $a_i>0$. Prove that, for any matrix $M$,
\[
 R(\theta,MX)=\|(M-I_d)\theta\|_2^2+\sigma^2\|M\|_F^2.
\]

\newpage
\item \textbf{(Upper bound: controlling the worst bias)} For a diagonal rule $\hat\theta_i=m_iX_i$, show that the maximum squared bias over $\Theta(A)$ is $\max_i(1-m_i)^2/a_i^2$. For a fixed $\lambda\ge0$, minimize the variance subject to this bias being at most $\lambda^2$. Show that the resulting weights are $m_i=(1-\lambda a_i)_+$, and deduce
\[
 R_L(\Theta(A))\le\inf_{\lambda\ge0}F(\lambda),\qquad
 F(\lambda)=\lambda^2+\sigma^2\sum_{i=1}^d(1-\lambda a_i)_+^2.
\]

\item \textbf{(Choosing the cutoff)} Show that $F$ has a unique minimizer $\lambda_*>0$, characterized by
\[
 \lambda_*=\sigma^2\sum_{i=1}^d a_i(1-\lambda_*a_i)_+.
\]
Why does this equation have exactly one positive solution? Let $w_i=(1-\lambda_*a_i)_+$ and $M_*=\operatorname{diag}(w_1,\ldots,w_d)$.

\item \textbf{(An exact semidefinite relaxation)} For any symmetric $D\succeq0$, prove that
\[
 \sup_{\theta^\top A\theta\le1}\theta^\top D\theta
 =\max_{\substack{\Sigma\succeq0\\\operatorname{tr}(A\Sigma)\le1}}\operatorname{tr}(D\Sigma)
 =\lambda_{\max}(A^{-1/2}DA^{-1/2}).
\]
Replacing $\theta\theta^\top$ by $\Sigma\succeq0$ drops a rank constraint. Why is there nevertheless no gap? Exhibit a rank-one maximizing $\Sigma$. \emph{Hint:} Set $H=A^{1/2}\Sigma A^{1/2}$ and use a top eigenvector.

\item \textbf{(Lower bound: a Gaussian Bayes problem)} Use E4 to express the worst-case risk of each fixed $MX$ as the maximum of its average risk under Gaussian priors $\mathcal N_d(0,\Sigma)$ satisfying $\operatorname{tr}(A\Sigma)\le1$. Then apply $\inf\sup\ge\sup\inf$ and B1 to show that
\[
 R_L(\Theta(A))\ge\max_{\substack{\Sigma\succeq0\\\operatorname{tr}(A\Sigma)\le1}}\operatorname{tr}\bigl[\sigma^2\Sigma(\Sigma+\sigma^2I_d)^{-1}\bigr].
\]
Explain why minimizing the average risk over linear rules gives the full Bayes risk here. Allow singular $\Sigma$, interpreting zero-variance directions as point masses.

\item \textbf{(Matching the two bounds)} Consider the diagonal covariance matrix
\[
 \Sigma_* =\operatorname{diag}(v_1^*,\ldots,v_d^*),\qquad
 v_i^*=\sigma^2\left(\frac1{\lambda_*a_i}-1\right)_+.
\]
Use E3 to show that $\operatorname{tr}(A\Sigma_*)=1$. Calculate its Bayes estimator and show that its Bayes risk is $\sigma^2\sum_iw_i=F(\lambda_*)$. Conclude that the lower bound in E5 matches the upper bound in E2. \emph{Hint:} Use $w_i^2+\lambda_*a_iw_i=w_i$.

\item \textbf{(Identifying the theorem and its scope)} State the exact value of $R_L(\Theta(A))$ and an estimator attaining it. Identify the result as Pinsker's linear minimax theorem and $M_*X$ as the Pinsker estimator. How does $a_i$ determine whether a coordinate is shrunk or discarded? The Gaussian prior in E6 is not supported on $\Theta(A)$. Explain why it still yields a valid lower bound for linear estimators, and why this argument does not establish minimaxity among all measurable estimators.
\end{enumerate}

\clearpage
\subsubsection*{(F) Application of Pinsker's theorem}
\phantomsection
\addcontentsline{toc}{subsubsection}{(F) Application of Pinsker's theorem}
\label{ex:pinsker-application}

For each $n\ge2$, observe the Gaussian sequence
\[
 X_i=\theta_i+\frac1{\sqrt n}Z_i,\qquad i=1,\ldots,d_n,
 \qquad Z_i\overset{\mathrm{i.i.d.}}{\sim}\mathcal N(0,1),
\]
where $d_n\ge\bigl\lceil\sqrt{(\log n)/2}\bigr\rceil$ and
\[
 \theta\in\Theta_n:=\left\{\theta\in\mathbb R^{d_n}:\sum_{i=1}^{d_n}e^{2i^2}\theta_i^2\le1\right\}.
\]
Under squared $\ell_2$ loss, define the minimax and linear minimax risks by
\[
 \begin{aligned}
 R_M(\Theta_n)&=\inf_{\hat\theta}\sup_{\theta\in\Theta_n}
   \mathbb E_\theta\|\hat\theta(X)-\theta\|_2^2,\\
 R_L(\Theta_n)&=\inf_{M\in\mathbb R^{d_n\times d_n}}\sup_{\theta\in\Theta_n}
   \mathbb E_\theta\|MX-\theta\|_2^2,
 \end{aligned}
\]
where the first infimum is over all measurable estimators.

\begin{enumerate}[label=\textcolor{AccentInk}{\sffamily\bfseries F\arabic*.},ref=F\arabic*,leftmargin=2.65em,labelsep=0.65em,itemsep=1em]
\item Using Pinsker's theorem, prove that, as $n\to\infty$,
\[
 R_M(\Theta_n)\sim R_L(\Theta_n)\sim\frac{\sqrt{\log n}}{\sqrt2\,n}.
\]
Here $\log$ denotes the natural logarithm and $u_n\sim v_n$ means $u_n/v_n\to1$.
\end{enumerate}
